\par\medskip
\hypertarget{native-079-article060}{}\pdfbookmark[2]{60.}{article-060}
\begin{center}60.\end{center}
Let $x$, $y$, $z$ denote the distances of the celestial body from three planes intersecting at right angles at the Sun, and $X$, $Y$, $Z$ the distances of the Earth (whether of its center or of a point on its surface) from the same planes. Clearly, $x-X$, $y-Y$, $z-Z$ will be the distances of the celestial body from three planes through the Earth parallel to those planes. These distances have the same relation to the distance of the body from the Earth and to its \textit{geocentric position}\footnote{In the broader sense. Strictly speaking, this expression refers to the case in which the line is drawn from the \textit{center} of the Earth.}, that is, the point on the celestial sphere onto which the line drawn from the Earth to the body is projected, as $x$, $y$, $z$ have to its distance from the Sun and its heliocentric position. Let $\Delta$ be the distance of the celestial body from the Earth. Imagine a perpendicular drawn on the celestial \SourceLeaf{080}sphere from the geocentric position to the great circle corresponding to the plane for the distances $z$. Let $a$ be the distance of the point of intersection from the positive pole of the great circle corresponding to the plane for the distances $x$. Finally, let $b$ be the length of the perpendicular, or the distance of the geocentric position from the great circle corresponding to the distances $z$. Then $b$ is the geocentric latitude or declination, according as the plane for the distances $z$ is the ecliptic or the equator. Similarly, $a+N$ is the geocentric longitude or right ascension if $N$ denotes the longitude of the pole of the plane for the distances $x$ in the former case, or its right ascension in the latter. We therefore have
\[
\begin{aligned}
x-X&=\Delta\cos b\cos a\\
y-Y&=\Delta\cos b\sin a\\
z-Z&=\Delta\sin b.
\end{aligned}
\]
The first two equations give $a$ and $\Delta\cos b$. The latter quantity (which must be taken positive), combined with the third equation, gives $b$ and $\Delta$.
